% GENERATED FILE. Edit canonical lessons, then run npm run generate:books.
% canonical-source-hash: fnv1a64:0e612f8aaf6ccdcc
% canonical-lessons: TE-C111-duvvena, TE-C111-sabbu, TE-C111-tuvvalu, TE-C111-kallajodu, TE-C111-gadiyaram

\chapter{Comb, Soap, Towel, Glasses, Clock}
\label{ch:te-comb-soap-towel-glasses-clock}
\hlchaptermodality{111}

\begin{chapteropening}
\textbf{By the end of this chapter:} \emph{I can say a comb, soap, a towel, glasses and a clock.}

Complete the last lesson of chapter 111: i can say a comb, soap, a towel, glasses and a clock.
\end{chapteropening}

\section[\texorpdfstring{\te{దువ్వెన} (duvvena)}{దువ్వెన (duvvena)}]{\te{దువ్వెన} (duvvena) — a comb}

\label{lesson:TE-C111-duvvena}



\begin{quote}
Before the new one: say the Telugu for a nail, then the Telugu for scissors.
\end{quote}

\subsection*{You'll want to know: \te{దువ్వెన}}
\textbf{\te{దువ్వెన}} — \emph{duvvena} — \textquotedblleft{}a comb\textquotedblright{}.

\begin{grammarlens}[title={What you've built}]
One more word for this chapter.
\end{grammarlens}

\subsection*{Your turn}
\begin{itemize}
\raggedright
  \item \emph{Say it:} \emph{duvvena}
  \item \emph{Say it:} \emph{duvvena}, once more
  \item \emph{Say it:} say \emph{duvvena}
  \item \emph{Recall:} say the Telugu for a nail, then the Telugu for scissors, then say \emph{duvvena} again
\end{itemize}

\begin{tcolorbox}[breakable,colback=teal!4,colframe=teal!35!black,title={Before you move on}]
What does \te{దువ్వెన} mean? (\textquotedblleft{}A comb\textquotedblright{}.) Say it once more.
\end{tcolorbox}
\section[\texorpdfstring{\te{సబ్బు} (sabbu)}{సబ్బు (sabbu)}]{\te{సబ్బు} (sabbu) — soap}

\label{lesson:TE-C111-sabbu}



\begin{quote}
Before the new one: say the Telugu for scissors, then the Telugu for a comb.
\end{quote}

\subsection*{You'll want to know: \te{సబ్బు}}
\textbf{\te{సబ్బు}} — \emph{sabbu} — \textquotedblleft{}soap\textquotedblright{}.

\begin{grammarlens}[title={What you've built}]
One more word for this chapter.
\end{grammarlens}

\subsection*{Your turn}
\begin{itemize}
\raggedright
  \item \emph{Say it:} \emph{sabbu}
  \item \emph{Say it:} \emph{sabbu}, once more
  \item \emph{Say it:} say \emph{sabbu}
  \item \emph{Recall:} say the Telugu for scissors, then the Telugu for a comb, then say \emph{sabbu} again
\end{itemize}

\begin{tcolorbox}[breakable,colback=teal!4,colframe=teal!35!black,title={Before you move on}]
What does \te{సబ్బు} mean? (\textquotedblleft{}Soap\textquotedblright{}.) Say it once more.
\end{tcolorbox}
\section[\texorpdfstring{\te{తువ్వాలు} (tuvvālu)}{తువ్వాలు (tuvvālu)}]{\te{తువ్వాలు} (tuvvālu) — a towel}

\label{lesson:TE-C111-tuvvalu}



\begin{quote}
Before the new one: say the Telugu for a comb, then the Telugu for soap.
\end{quote}

\subsection*{You'll want to know: \te{తువ్వాలు}}
\textbf{\te{తువ్వాలు}} — \emph{tuvvālu} — \textquotedblleft{}a towel\textquotedblright{}.

\begin{grammarlens}[title={What you've built}]
One more word for this chapter.
\end{grammarlens}

\subsection*{Your turn}
\begin{itemize}
\raggedright
  \item \emph{Say it:} \emph{tuvvālu}
  \item \emph{Say it:} \emph{tuvvālu}, once more
  \item \emph{Say it:} say \emph{tuvvālu}
  \item \emph{Recall:} say the Telugu for a comb, then the Telugu for soap, then say \emph{tuvvālu} again
\end{itemize}

\begin{tcolorbox}[breakable,colback=teal!4,colframe=teal!35!black,title={Before you move on}]
What does \te{తువ్వాలు} mean? (\textquotedblleft{}A towel\textquotedblright{}.) Say it once more.
\end{tcolorbox}
\section[\texorpdfstring{\te{కళ్ళజోడు} (kaḷḷajōḍu)}{కళ్ళజోడు (kaḷḷajōḍu)}]{\te{కళ్ళజోడు} (kaḷḷajōḍu) — glasses, spectacles}

\label{lesson:TE-C111-kallajodu}



\begin{quote}
Before the new one: say the Telugu for soap, then the Telugu for a towel.
\end{quote}

\subsection*{You'll want to know: \te{కళ్ళజోడు}}
\textbf{\te{కళ్ళజోడు}} — \emph{kaḷḷajōḍu} — \textquotedblleft{}glasses, spectacles\textquotedblright{}.

\begin{grammarlens}[title={What you've built}]
One more word for this chapter.
\end{grammarlens}

\subsection*{Your turn}
\begin{itemize}
\raggedright
  \item \emph{Say it:} \emph{kaḷḷajōḍu}
  \item \emph{Say it:} \emph{kaḷḷajōḍu}, once more
  \item \emph{Say it:} say \emph{kaḷḷajōḍu}
  \item \emph{Recall:} say the Telugu for soap, then the Telugu for a towel, then say \emph{kaḷḷajōḍu} again
\end{itemize}

\begin{tcolorbox}[breakable,colback=teal!4,colframe=teal!35!black,title={Before you move on}]
What does \te{కళ్ళజోడు} mean? (\textquotedblleft{}Glasses, spectacles\textquotedblright{}.) Say it once more.
\end{tcolorbox}
\section[\texorpdfstring{\te{గడియారం} (gaḍiyāraṁ)}{గడియారం (gaḍiyāraṁ)}]{\te{గడియారం} (gaḍiyāraṁ) — a clock, a watch}

\label{lesson:TE-C111-gadiyaram}



\begin{quote}
Before the new one: say the Telugu for a towel, then the Telugu for glasses.

Name the first Telugu month, \textbf{\te{చైత్రం}}, and the last, \textbf{\te{ఫాల్గుణం}}. Then say \textbf{\te{ఆదివారం}}, the day of the beginning.
\end{quote}

\subsection*{You'll want to know: \te{గడియారం}}
\textbf{\te{గడియారం}} — \emph{gaḍiyāraṁ} — \textquotedblleft{}a clock, a watch\textquotedblright{}.

\begin{grammarlens}[title={What you've built}]
One more word for this chapter.
\end{grammarlens}

\subsection*{Your turn}
\begin{itemize}
\raggedright
  \item \emph{Say it:} \emph{gaḍiyāraṁ}
  \item \emph{Say it:} \emph{gaḍiyāraṁ}, once more
  \item \emph{Say it:} say \emph{gaḍiyāraṁ}
  \item \emph{Recall:} say the Telugu for a towel, then the Telugu for glasses, then say \emph{gaḍiyāraṁ} again
\end{itemize}

\begin{tcolorbox}[breakable,colback=teal!4,colframe=teal!35!black,title={Before you move on}]
What does \te{గడియారం} mean? (\textquotedblleft{}A clock, a watch\textquotedblright{}.) Say it once more.
\end{tcolorbox}
